\documentclass[../../../include/open-logic-section]{subfiles}

\begin{document}
\olfileid{sth}{card-arithmetic}{simp}

\olsection{تبسيط الجمع والضرب}

يتبين أن جمع الكاردينالات وضربها فيما فوق المتناهي أمر
\emph{سهل للغاية}. وهذا ناشئ عن كون الكاردينالات أعداداً ترتيبية
(معينة)، ومن ثم فهي مرتبة ترتيباً جيداً ويمكن معالجتها بطريقة معينة.
غير أن إثبات ذلك \emph{ليس} سهلاً للغاية. ونحتاج، بدايةً، إلى تعريف
دقيق:

\begin{defn}
نعرّف \emph{ترتيباً معيارياً}، $\canonord$، على أزواج الأعداد
الترتيبية، فنشترط أن
$\tuple{\alpha_1, \alpha_2} \canonord
\tuple{\beta_1, \beta_2}$ إذا وفقط إذا تحقق أحد الشروط الآتية:
\begin{enumerate}
	\item $\max(\alpha_1, \alpha_2) < \max(\beta_1, \beta_2)$؛ أو
	\item $\max(\alpha_1, \alpha_2) = \max(\beta_1, \beta_2)$ و
	$\alpha_1 < \beta_1$؛ أو
	\item $\max(\alpha_1, \alpha_2) = \max(\beta_1, \beta_2)$ و
	$\alpha_1 = \beta_1$ و$\alpha_2 < \beta_2$.
\end{enumerate}
\end{defn}

\begin{lem}
$\tuple{\alpha \times \alpha, \canonord}$ ترتيب جيد، لكل عدد ترتيبي
$\alpha$.
\end{lem}

\begin{proof}
من الواضح أن $\canonord$ تام على $\alpha \times \alpha$. فلنفترض أن
لا $\tuple{\alpha_1, \alpha_2}$ أصغر من
$\tuple{\beta_1, \beta_2}$ بحسب $\canonord$، ولا الثاني أصغر من الأول.
عندئذ
$\max(\alpha_1, \alpha_2) = \max(\beta_1, \beta_2)$ و
$\alpha_1 = \beta_1$ و$\alpha_2 = \beta_2$، ومن ثم
$\tuple{\alpha_1, \alpha_2} = \tuple{\beta_1, \beta_2}$.

ولإثبات أنه ترتيب جيد، لتكن $X \subseteq \alpha\times\alpha$ غير
خالية. بما أن $\alpha$ عدد ترتيبي، يوجد $\delta$ هو أصغر عنصر في
$\Setabs{\max(\gamma_1, \gamma_2)}{\tuple{\gamma_1,
\gamma_2} \in X}$. احذف الآن من
$\Setabs{\tuple{\gamma_1,\gamma_2} \in X}{\max(\gamma_1, \gamma_2) =
\delta}$ كل الأزواج عدا التي لها أصغر إحداثي أول؛ ومن بين الباقي يكون
الزوج ذو أصغر إحداثي ثان هو أصغر عناصر $X$ بحسب $\canonord$.
\end{proof}
\noindent
والآن إلى ملاحظة صغيرة وبسيطة:

\begin{prop}\ollabel{simplecardproduct}
إذا كان $\cardeq{\alpha}{\beta}$، فإن
$\cardeq{\alpha \times \alpha}{\beta \times \beta}$.
\end{prop}

\begin{proof}
يكفي أن نجعل $f \colon \alpha \to \beta$ تستلزم الخريطة
$\tuple{\gamma_1, \gamma_2} \mapsto
\tuple{f(\gamma_1), f(\gamma_2)}$.
\end{proof}
\noindent
وسنضع كل هذا الآن موضع الاستعمال لإثبات لمّة حاسمة:
\begin{lem}\ollabel{alphatimesalpha}
لكل عدد ترتيبي غير متناهٍ $\alpha$،
$\cardeq{\alpha}{\alpha \times \alpha}$.
\end{lem}

\begin{proof}
للبرهان بالخلف، ليكن $\alpha$ أصغر عدد ترتيبي غير متناهٍ لا تصح له هذه
النتيجة. تبين \olref[sfr][siz][zigzag]{natsquaredenumerable} أن
$\cardeq{\omega}{\omega\times\omega}$، ولذلك
$\omega \in \alpha$. فضلاً عن ذلك، $\alpha$ كاردينال: ولنفترض العكس
للبرهان بالخلف؛ عندئذ $\card{\alpha} \in \alpha$، فيكون
$\cardeq{\card{\alpha}}{\card{\alpha} \times \card{\alpha}}$ بحسب
الفرض، و$\cardeq{\card{\alpha}}{\alpha}$ بحسب التعريف؛ ومن ثم
$\cardeq{\alpha}{\alpha\times\alpha}$ بحسب
\olref{simplecardproduct}.

والآن، لكل $\tuple{\gamma_1, \gamma_2} \in \alpha \times \alpha$،
انظر إلى القطعة:
\begin{align*}
	\text{Seg}(\gamma_1, \gamma_2) &= \Setabs{\tuple{\delta_1, \delta_2} \in \alpha \times \alpha}{\tuple{\delta_1, \delta_2} \canonord \tuple{\gamma_1, \gamma_2}}
\end{align*}
ضع $\gamma = \max(\gamma_1, \gamma_2)$، ولاحظ أن
$\tuple{\gamma_1, \gamma_2} \canonord \tuple{\gamma+1, \gamma + 1}$.
إذا كان $\gamma$ متناهياً، فالقطعة متناهية، ولذلك كارديناليتها أصغر من
الكاردينال غير المتناهي $\alpha$. أما إذا كان $\gamma$ غير متناهٍ،
فنحصل على:
\begin{align*}
	\text{Seg}(\gamma_1, \gamma_2) &
	\precsim ((\gamma \ordplus 1)\ordtimes (\gamma \ordplus 1))\\
	&\approx (\gamma \ordtimes \gamma)
	\text{، بحسب \olref[ord-arithmetic][using-addition]{ordinfinitycharacter} و
	\olref{simplecardproduct}}\\
	&\approx \gamma \text{، بحسب فرض الاستقراء}\\
	& \prec \alpha\text{، لأن $\alpha$ كاردينال.}
\end{align*}
إذن $\ordtype{\alpha\times \alpha, \canonord} \leq \alpha$، ومن ثم
$\cardle{\alpha \times \alpha}{\alpha}$. وبما أن
$\cardle{\alpha}{\alpha \times \alpha}$ بالطبع، فالنتيجة تتبع من
مبرهنة شرودر--برنشتاين.
\end{proof}

وأخيراً نصل إلى النتيجة المبسِّطة:

\begin{thm}\ollabel{cardplustimesmax}
إذا كان $\cardfont{a}, \cardfont{b}$ كاردينالين غير متناهيين، فإن:
\[
	\cardfont{a}
\cardtimes \cardfont{b} = \cardfont{a} \cardplus \cardfont{b} =
\text{max}(\cardfont{a}, \cardfont{b}).
\]
\end{thm}

\begin{proof}
من دون إخلال بالعموم، افترض أن
$\cardfont{a} = \max(\cardfont{a}, \cardfont{b})$. وباستعمال
\olref{alphatimesalpha} نحصل على
$\cardfont{a}\cardtimes\cardfont{a} = \cardfont{a} \leq \cardfont{a}
\cardplus \cardfont{b} \leq \cardfont{a} \cardplus \cardfont{a} \leq
\cardfont{a} \cardtimes \cardfont{a}$.
\end{proof}\noindent
وبالمثل، إذا كان $\cardfont{a}$ غير متناهٍ، فإن اتحاد عدد
$\cardfont{a}$ من المجموعات التي لا يزيد حجم كل منها على
$\cardfont{a}$ لا يزيد حجمه على $\cardfont{a}$:

\begin{prop}\ollabel{kappaunionkappasize}
ليكن $\cardfont{a}$ كاردينالاً غير متناهٍ. ولكل عدد ترتيبي
$\beta \in \cardfont{a}$، لتكن $X_\beta$ مجموعة تحقق
$\card{X_\beta} \leq \cardfont{a}$. عندئذ
$\card{\bigcup_{\beta \in \cardfont{a}} X_\beta} \leq \cardfont{a}$.
\end{prop}

\begin{proof}
لكل $\beta \in \cardfont{a}$، ثبّت !!a{injection}
$f_\beta \colon X_\beta \to \cardfont{a}$.\footnote{كيف «ثُبِّتت» هذه
الدوال؟ انظر \olref[sth][choice][countablechoice]{sec}.} عرّف
!!a{injection} $g \colon
\bigcup_{\beta \in \cardfont{a}} X_\beta \to \cardfont{a} \times
\cardfont{a}$ بوضع $g(v) = \tuple{\beta, f_\beta(v)}$، حيث
$v \in X_\beta$ ولا يكون $v \in X_\gamma$ لأي $\gamma \in \beta$.
والآن، بحسب \olref{cardplustimesmax}،
$\bigcup_{\beta \in \cardfont{a}} X_\beta \preceq
\cardfont{a} \times \cardfont{a} \approx \cardfont{a}$.
\end{proof}

\end{document}
